\section{Effectivity with a prescribed correction support}
\label{sec:effectivity}

The signed approximation separates two kinds of data: the negative
part is confined to the permitted correction support, while the ample
part can absorb a small multiple of any prescribed rational divisor.
We use that ample part to insert the pseudoeffective divisor of
Theorem~\ref{thm:controlled-boundary}.  The resulting boundary is
effective, big and klt, so big-boundary nonvanishing applies to its
adjoint.  We first record the passage from its real-linear conclusion
to the rational linear equivalence required here.

\begin{lemma}[Rational recovery]\label{lem:rational-recovery}
Let $D$ be a rational divisor on a normal projective variety.
If $D$ is real-linearly equivalent to an effective real divisor, then
it is rationally linearly equivalent to an effective rational divisor.
\end{lemma}
\begin{proof}
Write an effective real representative as
$D+\sum_{j=1}^s a_j\operatorname{div}(f_j)$, with $a_j\in\mathbb R$
and rational functions $f_j$.  The supports of $D$ and of these finitely
many principal divisors give a finite system of rational linear
inequalities asserting nonnegativity of every coefficient.
The known real solution makes the resulting rational polyhedron
nonempty.  Such a polyhedron has a rational point: take a point in the
relative interior of its smallest supporting face, solve that face's
rational linear equations over $\mathbb Q$, and approximate within
its rational affine space while preserving the remaining strict
inequalities.  Substituting that rational point gives the assertion.
\end{proof}

\begin{proof}[Proof of Theorem~\ref{thm:controlled-boundary}]
Fix a small ample rational divisor $A$ whose smooth metric can be
subtracted from the given metric on $-K_Y+C$ while retaining its global
strict curvature lower bound.  Apply
Theorem~\ref{thm:weighted-approximation} with $J=C$ and choose $r>1$.
On the resulting resolution $\sigma:Y'\to Y$, write
$\Psi=\Psi^+-\Psi^-$ and retain the free system $|\mathrm{Mov}|$.
The signed adjoint identity is
\[
 K_{Y'}+\Psi^++\mathrm{Mov}/r+\sigma^*A
       \sim_{\mathbb Q}\Psi^-.
\]
Choose a general $D_{\mathrm{mov}}\in|\mathrm{Mov}|$ meeting the
resolved support transversely and without common components.
All these data are fixed before choosing the pseudoeffective divisor $M$.

The ample cone is open, so $A+tM$ is ample for a sufficiently small
positive rational $t$.  Choose an effective rational representative
$U_t\sim_{\mathbb Q}A+tM$ by dividing a general member of a sufficiently
high very ample multiple.  Bertini permits $\sigma^*U_t$ and $D_{\mathrm{mov}}$ to meet
the resolved support with simple normal crossings, with no common
components; the new coefficients are less than one.  Thus
\begin{equation}\label{eq:effective-big-boundary}
 \Delta_t=\Psi^++D_{\mathrm{mov}}/r+\sigma^*U_t
\end{equation}
is effective and klt.  Its class is big because it contains the
effective summand $\sigma^*U_t$, the pullback of an ample rational class.
The adjoint identity becomes
\begin{equation}\label{eq:effectivity-adjoint}
 K_{Y'}+\Delta_t\sim_{\mathbb Q}\Psi^-+t\sigma^*M.
\end{equation}
Its right-hand side is pseudoeffective.

Big-boundary nonvanishing of Birkar--Cascini--Hacon--McKernan
\cite[Theorem D]{BCHM2010} therefore gives an effective real divisor
real-linearly equivalent to $K_{Y'}+\Delta_t$.
Lemma~\ref{lem:rational-recovery} makes this a rational linear
equivalence with an effective rational divisor.  Finally,
$F=\sigma_*\Psi^-$ is effective and supported on $C$ by
Lemma~\ref{lem:weighted-resolution}.  Push
\eqref{eq:effectivity-adjoint} down to $Y$ and divide by $t$.
The theorem follows with $C_0=F/t$.
\end{proof}

\begin{proposition}[A fixed support correction before scaling]
\label{prop:ample-perturbation}
Under the hypotheses of Theorem~\ref{thm:controlled-boundary}, there
is an effective rational divisor $F$ supported on $C$ such that for
every pseudoeffective rational divisor $M$, some positive rational
$t$ makes $F+tM$ rationally linearly equivalent to an effective
rational divisor.
\end{proposition}
\begin{proof}
In the preceding proof, the approximation and hence
$F=\sigma_*\Psi^-$ were chosen independently of $M$.  Only $t$ and the
general ample representative depended on $M$.
\end{proof}

When $C=0$, the correction vanishes.  For nonzero $C$, it is the
integrability of $s_C$ that controls the support of $\Psi^-$; no
unweighted frame-integrability hypothesis has been imposed.
The divisor to which nonvanishing is applied is the effective boundary
\eqref{eq:effective-big-boundary}.

\subsection{Preserving integrability under a pseudoeffective perturbation}

One can also insert $M$ at the level of metrics, before approximation.
This gives analytic information in addition to the effectivity proved
above: a small pseudoeffective perturbation preserves the integrability
of the distinguished section and the global strict curvature bound.

\begin{lemma}\label{lem:weighted-pseudoeffective-perturbation}
Under the hypotheses of Theorem~\ref{thm:controlled-boundary}, for
every pseudoeffective rational divisor $M$ there is a positive rational
$\epsilon$ such that $-K_Y+C+\epsilon M$ has a metric with a global
strict curvature lower bound and multiplier ideal containing
$\mathcal O_Y(-C)$.
\end{lemma}
\begin{proof}
Write $\phi$ for the given metric and choose a semipositive singular
metric with weights $\theta$ on $M$.  Strong openness
\cite{GuanZhou2015} gives a common $\tau>0$ on a finite relatively
compact coordinate covering such that
$|s_C|^2e^{-(1+\tau)\phi}$ is locally integrable.  Put
$q=(1+\tau)/\tau$.  H\"older's inequality with respect to the measure
$|s_C|^2d\lambda$ gives
\[
 \int |s_C|^2e^{-\phi-\epsilon\theta}
 \le
 \left(\int |s_C|^2e^{-(1+\tau)\phi}\right)^{1/(1+\tau)}
 \left(\int |s_C|^2e^{-q\epsilon\theta}\right)^{1/q}.
\]
Small-exponent integrability of plurisubharmonic functions
\cite{Skoda1972} makes the second factor finite for sufficiently small
positive rational $\epsilon$.  The finite covering permits one choice
of $\epsilon$.  The weight $\phi+\epsilon\theta$ retains the global
strict curvature lower bound because the metric on $M$ is semipositive.
\end{proof}

Subtract a sufficiently small ample rational $A$ from the perturbed
metric, retaining its strict lower bound.  Applying
Lemma~\ref{lem:weighted-resolution} to $-K_Y+C+\epsilon M-A$ with
$r>1$ then gives a second boundary construction.
For its signed divisor $\Psi$ and a general free member
$D_{\mathrm{mov}}$, add a general small-coefficient representative
$U\sim_{\mathbb Q}A$.  Then
$\Delta=\Psi^++D_{\mathrm{mov}}/r+\sigma^*U$ is effective, big and klt,
and $K_{Y'}+\Delta\sim_{\mathbb Q}\epsilon\sigma^*M+\Psi^-$.
Nonvanishing, rational recovery and pushforward apply as before.
Here the approximation itself depends on the metric on $M$;
the support conclusion is unchanged.
